\section{The question, the results, and their provenance}
\label{sec:context}

Throughout the paper $m$ is Lebesgue measure and $\N=\{1,2,\ldots\}$.
For a measurable $1$-periodic set $H$ write
\[
 \rho(H)=m(H\cap[0,1)).
\]
A nontrivial affine copy of a set $A\subset\R$ has the form $x+sA$,
where $x\in\R$ and $s\ne0$. Put $G_q=\{q^n:n\ge1\}$ for $0<q<1$.
The distinction between one fixed $q$ and the entire continuum of ratios
is central to the results.

\subsection{A published question with different quantifiers}

The Erd\H{o}s similarity conjecture asks whether every infinite subset of
$\R$ has a positive-measure set avoiding all its nontrivial affine copies.
Burgin, Goldberg, Keleti, MacMahon, and Wang posed the following separate
question in \cite[Question~1]{BGKMW}: must every measurable set of positive
measure contain $\{x+sq^n:n\ge1\}$ for \emph{some} triple $(x,s,q)$ with
$s\ne0$ and $0<q<1$? A negative answer requires a single set that excludes
all three parameters at once.

Theorem~1.1 of their paper constructs a compact set having zero as a density
point but containing no unshifted progression $\{sq^n:n\ge1\}$. That result
does not exclude progressions converging to another center. Their
almost-everywhere avoidance statement also leaves exceptional parameter
values. The survey \cite{JLM} provides the historical context; it predates
the October 2026 manuscript used below.

The OpenAI manuscript \cite[Theorem~1.1]{OpenAI84} proves that for each
\emph{fixed} $q$ there are compact subsets of $[0,1]$ of measure arbitrarily
close to one avoiding all affine copies of $G_q$. Its introduction explicitly
permits the avoiding set to depend on $q$. Summable intersections therefore
handle any specified countable collection of ratios. They do not by
themselves handle every real ratio. This paper supplies that missing
uniform argument.

\begin{theorem}[Uniform geometric avoidance]\label{thm:main}
For every $0<\eps<1$ there exists a closed, symmetric, $1$-periodic set
$F_\eps\subset\R$ such that
\begin{equation}\label{eq:main-density}
 m(F_\eps\cap[u,u+1])\ge1-\frac\eps2>1-\eps
 \qquad(u\in\R),
\end{equation}
and, for every $x\in\R$, $s\in\R\setminus\{0\}$, and $q\in(0,1)$,
\begin{equation}\label{eq:main-avoidance}
 x+sq^n\notin F_\eps\quad\text{for infinitely many }n\ge1.
\end{equation}
Consequently $K_\eps=F_\eps\cap[0,1]$ is compact, has measure greater
than $1-\eps$, and answers \cite[Question~1]{BGKMW} negatively.
\end{theorem}

The conclusion is stronger than avoidance at a prescribed limit point.
Every point, every signed dilation, and every contracting ratio are covered
by the same set. Periodicity in \cref{thm:main} is possible because passing
to a tail of a geometric progression reduces its scale. The construction
only needs dyadic contractions of periodic sets.

\begin{theorem}[All finite positive exponential mixtures]\label{thm:mixtures}
For every $0<\eps<1$ there is a closed, symmetric, $1$-periodic set
$F^*_\eps\subset\R$ satisfying \eqref{eq:main-density} and the following
property. For every finite $d\ge1$, every $q_1,\ldots,q_d\in(0,1)$,
every $c_1,\ldots,c_d\ge0$ with $\sum_i c_i>0$, every $x\in\R$, and
every $s\ne0$, infinitely many terms of
\begin{equation}\label{eq:mixture-orbit}
 x+s\sum_{i=1}^d c_iq_i^n,\qquad n\ge1,
\end{equation}
lie outside $F^*_\eps$. All choices, including the number $d$ of
components, are simultaneous.
\end{theorem}

The terms in \eqref{eq:mixture-orbit} are synchronized sums: the exponent
$n$ is the same in every summand. This is a different pattern from the
sumset in which all component exponents vary independently. Repeated
bases and zero coefficients are allowed. The proof is an application of
a general theorem for finite-dimensional polynomial families in
\cref{sec:algebraic}. That section also proves a stronger extension to
all nonzero real exponential polynomials with positive real bases,
allowing signed coefficients, and deduces the corresponding theorem
for homogeneous constant-coefficient real recurrences with roots in $(0,1)$.

\begin{theorem}[Effective compact sets and finite margins]\label{thm:effective-intro}
For rational $0<\eps<1$, the compact sets in \cref{thm:main,thm:mixtures}
can be chosen effectively closed, uniformly in $\eps$, with computable
Lebesgue measure. More precisely, they have nested approximations $K_N$
by finite unions of closed rational intervals such that
\begin{equation}\label{eq:effective-rate-intro}
 K=\bigcap_{N\ge1}K_N,
 \qquad 0\le m(K_N)-m(K)\le\eps\,2^{-N}.
\end{equation}
For any rational compact parameter box contained in the admissible domain,
with the dilation bounded away from zero, one can algorithmically obtain a finite index bound and a
positive rational avoidance margin. The precise certificate is given in
\cref{thm:certificates}.
\end{theorem}

Effectively closed here means that the open complement can be enumerated
as a union of rational open intervals. It does not assert decidable
membership, or computability in the Hausdorff representation of compact
sets. The finite search procedures proved to terminate below are not
claimed to be efficient.

\subsection{What was reviewed and what is new}

The repository review used the revisions recorded in \cref{sec:provenance}.
The OpenAI catalogue contains manuscripts at different verification stages;
its own README makes that distinction \cite{OpenAICatalog}. We inspected
the actual source of the geometric similarity manuscript, not just its
catalogue description. In ProveIt we reviewed the repository overview,
Ramsey refinement status, and selected matching and analysis material
\cite{ProveIt}. The large collection of existing Gowers refinements made
another minor local refinement a less attractive target. The present
question offers a precise published target and a useful connection to
ProveIt's work on exact certificates, polynomial decision procedures, and
exponential sequences.

\begin{center}
\begin{tabular}{@{}>{\raggedright\arraybackslash}p{0.30\textwidth}>{\raggedright\arraybackslash}p{0.62\textwidth}@{}}
\toprule
Ingredient & Status in this article\\
\midrule
Ordered-tree routing; selector and terminal tables & Adapted, with full
attribution, from the fixed-ratio proof in \cite{OpenAI84}. The probability
argument is reproduced in our notation and checked directly.\\[3pt]
Common-scale subsequence for varying $q$ & New uniformization step:
$\nu_j(q)=\min\{n:q^n\le Q^j\}$. The random set is independent of $q$.\\[3pt]
Two-parameter cell count & New use of lines in logarithmic coordinates,
including jump points, edges, and vertices.\\[3pt]
Polynomial-family extension & New application of the classical
sign-condition bound \cite{BPR}; yields \cref{thm:mixtures}.\\[3pt]
Periodic global assembly and effective certificates & Developed here;
uses tail closure, summable budgets, and real quantifier elimination.\\
\bottomrule
\end{tabular}
\end{center}

The principal theorem is proved with finite probability and elementary
measure theory. The general polynomial-family extension additionally uses
a stated classical theorem of real algebraic geometry. The effective
results use decidability of real polynomial inequalities. No flagship
claim elsewhere in either repository is an unproved hypothesis of these
arguments.

\subsection{Road map}

We first synchronize all ratios in a compact interval to one family of
spatial scales. A finite tree gives independent opportunities for a hit
while keeping the mean measure small. The first default choice in the
route of a fixed center identifies the relevant tests. A line arrangement
then replaces the continuum of ratio--dilation parameters by finitely many
representatives. Compactness repairs the remaining centers. Finally,
countable summable assembly covers all ratios and scales. The same scheme
with polynomial sign conditions handles exponential mixtures.
